\section{Problem Tanimi ve Kapsam}

Bu proje, EGO benzeri bir toplu tasima takip sisteminin teknik cekirdeginin sentetik veri ile modellenmesini hedeflemektedir. Temel problem, hareket halindeki otobuslerden gelen telemetri verisinin gercek zamanli olarak toplanmasi, bulut uzerinde islenmesi, operasyonel olarak anlamli hale getirilmesi ve son kullaniciya anlik olarak sunulmasidir. Bu tanim, calismanin odagini basit bir harita gosteriminden ayirmakta; veri akisi, bulut servis entegrasyonu ve olay tabanli isleme mantigini merkeze almaktadir.

Projede ham veri ile turetilen veri birbirinden acik bicimde ayrilmistir. Ham veri, simulator tarafinda uretilen \texttt{bus\_id}, \texttt{line\_id}, \texttt{lat}, \texttt{lon}, \texttt{speed\_kmh}, \texttt{next\_stop\_id}, \texttt{next\_stop\_name}, \texttt{boarding\_count} ve \texttt{timestamp} alanlarindan olusmaktadir. Buna karsilik \texttt{estimated\_eta\_sec}, \texttt{estimated\_alighting\_count}, \texttt{estimated\_occupancy\_score}, \texttt{estimated\_occupancy\_level} ve \texttt{is\_delayed} alanlari bulut tarafinda hesaplanmistir. Bu ayrim, sistemin yalnizca veri gosteren bir simulasyon olmadigini, veriyi yorumlayan ve zenginlestiren bir bulut isleme katmani kurdugunu gostermektedir.

Calisma kapsaminda sistem, MVP mantigiyla sinirlandirilmis ve uc hat, hat basina birden fazla otobus ve tanimli durak/rota dizileri uzerinden ilerleyecek sekilde modellenmistir. Gercek EGO altyapisina baglanti kurulmamistir; fiziksel otobus cihazlari, GPS modulleri veya kart okuyucular kullanilmamistir. Makine ogrenmesi tabanli tahminleme, kullanici girisi, mobil istemci, dinamik rota optimizasyonu ve gercek trafik verisi entegrasyonu gibi genisletilebilir bilesenler bu asamada bilerek kapsam disinda birakilmistir. Boylece proje kapsami, ders gereksinimlerine uygun olarak gercek zamanli veri akisi, bulut isleme, depolama ve gorsellestirme ekseninde kontrollu bicimde tutulmustur.

Bu kapsam siniri, raporun degerlendirme mantigi acisindan da onemlidir. Basari olcutu, gercek bir belediye sisteminin tum karmasikligini taklit etmek degil; secilen mimarinin veri uretimi, MQTT tabanli iletim, AWS uzerinde zenginlestirme, DynamoDB ile depolama ve dashboard ile canli izleme asamalarini uctan uca yerine getirebilmesidir.
